\section{Results}

% ===================================================================
\subsection{Study 1: EmotiView Pilot Study}
\label{sec:results_ev1}
% ===================================================================

\noindent\textit{Note: EV1 served as a proof-of-concept testbed; results are reported in compact tabular form rather than as per-participant plots. Full bootstrap distributions, spectral plots, and fNIRS hemodynamic traces are available in the interactive HTML result archive (\texttt{EV\_results.html}). Inferential statistics are consolidated in Table~\ref{tab:pilot_anova}.}

\subsubsection{Participant Characteristics and Data Quality}

Five participants (3 females, 2 males; mean age = 24.2 years, SD = 2.4, range = 21-27) completed the full experimental protocol. All participants were right-handed (Edinburgh Handedness Inventory laterality quotient > +40), native German speakers, and reported no history of neurological or psychiatric disorders. Data quality assessment revealed successful multimodal recording across all participants, with \gls{LSL} synchronization maintaining temporal alignment across \gls{fNIRS}, \gls{EEG}, \gls{ECG}, and \gls{EDA} streams. Automated pipeline processing completed successfully for all five participants, generating 11 terminal output modules per participant (questionnaire assessments, physiological metrics, and neurophysiological measures). Full per-participant metrics are embedded in the PDF/A-compliant pipeline output documents.

\subsubsection{Baseline Trait Measures}

Baseline questionnaires assessed trait motivational tendencies and pre-experimental affective state. \gls{BIS}/\gls{BAS} scores revealed individual differences in behavioral inhibition and activation systems. Mean \gls{BIS} score across participants was M = 18.4 (SD = 3.2), indicating moderate sensitivity to punishment cues. \gls{BAS} subscales showed: Drive M = 10.2 (SD = 2.1), Fun-Seeking M = 11.6 (SD = 1.8), and Reward Responsiveness M = 16.8 (SD = 2.4). These values fall within normative ranges for young adult populations \parencite{strobelDeutschsprachigeVersionBIS2006}, suggesting the sample exhibited typical motivational profiles. Pre-experimental \gls{PANAS} indicated participants began the session in neutral-to-positive affective states (Positive Affect mean = 28.6, SD = 4.2; Negative Affect mean = 12.4, SD = 2.8), consistent with typical laboratory baseline states \parencite{watsonDevelopmentValidationBrief1988}. Per-participant bar charts for both instruments are available in the interactive HTML result archive (\texttt{EV\_results.html}).

\subsubsection{Emotion Induction Validation}

Post-stimulus \gls{SAM} ratings confirmed successful emotion induction across valence conditions (Table~\ref{tab:sam}). Participants rated emotional videos on 9-point scales for valence (1 = very negative, 9 = very positive) and arousal (1 = very calm, 9 = very aroused). Within-participant patterns showed systematic differentiation consistent with the circumplex model of affect \parencite{russellCircumplexModelAffect1980}: positive videos elicited high valence/moderate arousal, negative videos low valence/high arousal, and neutral videos mid-range valence/low arousal.

\begin{table}[htbp]
    \centering
    \small
    \caption{\gls{SAM} Valence and Arousal Ratings by Condition (Group Mean $\pm$ SD, $N = 5$)}
    \label{tab:sam}
    \begin{tabular}{lrrr}
        \hline
        \textbf{Dimension} & \textbf{Negative} & \textbf{Neutral} & \textbf{Positive} \\
        \hline
        Valence (1--9) & $2.4 \pm 1.1$ & $5.2 \pm 0.8$ & $6.8 \pm 1.2$ \\
        Arousal (1--9) & $6.2 \pm 1.6$ & $3.1 \pm 0.9$ & $5.6 \pm 1.4$ \\
        \hline
    \end{tabular}
    \par\vspace{4pt}
    \noindent\parbox{\linewidth}{\scriptsize 9-point scales (valence: 1 = very negative, 9 = very positive; arousal: 1 = very calm, 9 = very aroused). Per-participant bar charts available in \texttt{EV\_results.html}.}
\end{table}

Basic emotion evaluations using a 7-category scheme (joy, surprise, fear, disgust, anger, sadness, neutral) revealed condition-specific emotion profiles: positive videos predominantly elicited joy, negative videos showed distributed patterns across aversive categories, and neutral videos were appropriately rated as neutral. Emotional adjective ratings (11-item scale) further confirmed granular affective differentiation across conditions \parencite{maffeiEMOVIEExperimentalMOVies2019}. Per-participant bar charts for both instruments are available in the interactive HTML result archive (\texttt{EV\_results.html}).

\subsubsection{Sympathetic Arousal: Electrodermal Activity}

\gls{EDA} analysis focused on phasic skin conductance changes reflecting sympathetic nervous system activation during emotional processing. Bootstrap resampling analysis (10,000 iterations, percentile method, $\alpha$ = 0.05) was applied to epoch-level amplitude data (peak-baseline conductance change in µS) to compute robust within-participant confidence intervals for each experimental condition.

Results revealed condition-dependent sympathetic responses with individual variability in magnitude (Table~\ref{tab:ev1_descriptive}). Participants EV\_002, EV\_003, and EV\_007 demonstrated non-zero \gls{EDA} amplitude across conditions. Bootstrap confidence intervals indicated stable within-participant amplitude estimates; however, no consistent directional pattern was observed across the NEG, NEU, and POS conditions. Within-participant one-way \gls{ANOVA} did not reach significance in any participant (all $F < 1.36$, $p > .33$; see Table~\ref{tab:pilot_anova}).

Bootstrap CI widths varied across participants and conditions, reflecting differences in trial-to-trial consistency. Narrow CIs indicated stable, replicable responses, while wider CIs reflected greater intra-individual variability in sympathetic reactivity.

Participants EV\_004 and EV\_008 showed minimal or absent \gls{EDA} responses. EV\_008 amplitude values ($\approx 2$--3\,\textmu S) were two orders of magnitude lower than EV\_002 through EV\_007 ($\approx 100$\,\textmu S); the source of this discrepancy (electrode contact, signal processing units, or genuine individual difference) is not resolved and EV\_008 \gls{EDA} data should be interpreted with caution \parencite{boucseinElectrodermalActivity2012}.

\subsubsection{Parasympathetic Regulation: Heart Rate Variability}

\gls{HRV} analysis quantified vagal tone via \gls{RMSSD} (root mean square of successive NN interval differences, measured in milliseconds), a time-domain metric sensitive to parasympathetic influence on cardiac regulation \parencite{malikHeartRateVariability1996, shafferOverviewHeartRate2017}. Bootstrap resampling was applied to epoch-level RMSSD values to assess within-participant condition effects with non-parametric confidence intervals.

\gls{RMSSD} values varied by condition across participants (Table~\ref{tab:ev1_descriptive}). Within-participant one-way \gls{ANOVA} reached significance in 4 of 5 participants (EV\_002: $F(2,189) = 27.80$, $p < .001$, $\eta^2 = .23$; EV\_003: $F(2,190) = 11.02$, $p < .001$, $\eta^2 = .10$; EV\_007: $F(2,190) = 3.87$, $p = .023$, $\eta^2 = .04$; EV\_008: $F(2,191) = 18.71$, $p < .001$, $\eta^2 = .16$; EV\_004: $F(2,170) = 2.49$, $p = .086$, $\eta^2 = .03$, ns). See Table~\ref{tab:pilot_anova} for the full per-participant breakdown.

\noindent\textbf{Data quality note.} \gls{RMSSD} values for EV\_002 ($M_{\text{NEG}} = 818$\,ms, $M_{\text{POS}} = 1540$\,ms), EV\_003 ($M = 199$--$281$\,ms), EV\_007 ($M = 98$--$133$\,ms), and EV\_008 ($M = 91$--$105$\,ms) substantially exceed the physiologically expected range for ECG-based \gls{RMSSD} in resting or mildly active adults ($\approx 20$--$50$\,ms \parencite{malikHeartRateVariability1996, shafferOverviewHeartRate2017}). Only EV\_004 ($M = 44$--$47$\,ms) falls within the expected range, and EV\_004 is the sole non-significant participant. The anomalous values most likely reflect R-peak detection errors (missed beats or ectopic-beat misclassification) in the automated processing pipeline; the raw ECG waveforms were not manually verified. The \gls{RMSSD} statistics for EV\_002, EV\_003, EV\_007, and EV\_008 should therefore be regarded as preliminary and their directional interpretation withheld pending signal-quality validation.

Bootstrap confidence intervals revealed individual differences in response stability. EV\_003 demonstrated particularly narrow \gls{RMSSD} bootstrap CIs across all conditions, though absolute values ($M = 199$--$281$\,ms) lie outside the physiologically expected range noted above. EV\_007 and EV\_008 exhibited wider CIs, indicating greater trial-to-trial variability in computed RMSSD values. Directional relationships between \gls{EDA} and \gls{HRV} bootstrap estimates across conditions are available in the interactive HTML result archive (\texttt{EV\_results.html}).

\begin{table}[htbp]
    \centering
    \small
    \setlength{\tabcolsep}{5pt}
    \caption{EV1 Peripheral Autonomic Descriptive Summary ($N = 5$)}
    \label{tab:ev1_descriptive}
    \begin{tabular}{lp{3.2cm}p{3.2cm}p{4.2cm}}
        \hline
        \textbf{Participant} & \textbf{EDA range ($\mu$S)} & \textbf{RMSSD range (ms)} & \textbf{Notes} \\
        \hline
        EV\_002 & $\approx 100$  & 818--1540 & RMSSD likely inflated by R-peak detection errors \\
        EV\_003 & $\approx 100$  & 199--281  & RMSSD likely inflated by R-peak detection errors \\
        EV\_004 & minimal        & 44--47    & Only physiologically plausible RMSSD estimate \\
        EV\_007 & $\approx 100$  & 98--133   & RMSSD likely inflated by R-peak detection errors \\
        EV\_008 & $\approx 2$--3 & 91--105   & EDA scale anomaly; RMSSD likely inflated \\
        \hline
    \end{tabular}
    \par\vspace{4pt}
    \noindent\parbox{\linewidth}{\scriptsize Ranges reflect within-participant variability across the three valence conditions (NEG, NEU, POS). EDA: peak-to-baseline conductance change. RMSSD: root mean square of successive NN differences. Inflated RMSSD values most likely reflect R-peak detection artefacts; see data quality note above. Condition-level bootstrap estimates available in \texttt{EV\_results.html}.}
\end{table}

\subsubsection{Cortical Oscillatory Dynamics: EEG Power Spectral Density}

\gls{EEG} power spectral analysis characterized condition-dependent modulation of cortical oscillations in alpha (8--13\,Hz), beta (13--30\,Hz), and theta (4--8\,Hz) frequency bands. Frontal alpha asymmetry indices (\gls{FAI}), computed as $\ln P_{\text{right}} - \ln P_{\text{left}}$ for electrode pair F4/F3, were examined for condition-dependent modulation.

Individual variability was observed in both absolute spectral power and \gls{FAI}. \gls{FAI} reached significance in EV\_003 ($F(2,181) = 4.50$, $p = .012$, $\eta^2 = .05$) and EV\_007 ($F(2,195) = 4.73$, $p = .010$, $\eta^2 = .05$), while EV\_002, EV\_004, and EV\_008 did not reach significance (all $p > .59$). Key per-participant effect sizes are summarised in Table~\ref{tab:ev1_eeg}; full \gls{ANOVA} statistics are in Table~\ref{tab:pilot_anova}. Per-participant spectral plots are available in the interactive HTML result archive (\texttt{EV\_results.html}).

\begin{table}[htbp]
    \centering
    \small
    \setlength{\tabcolsep}{5pt}
    \caption{EV1 EEG Spectral Summary: Selected Significant Effects per Participant}
    \label{tab:ev1_eeg}
    \begin{tabular}{lp{3.5cm}rrl}
        \hline
        \textbf{Participant} & \textbf{Measure} & $\boldsymbol{F}$ & $\boldsymbol{\eta^2}$ & \textbf{sig} \\
        \hline
        EV\_002 & Frontal theta   & 4.00  & .048 & * \\
        EV\_003 & Frontal theta   & 34.92 & .297 & *** \\
        EV\_003 & \gls{FAI}       & 4.50  & .047 & * \\
        EV\_007 & Parietal alpha  & 3.13  & .031 & * \\
        EV\_007 & Parietal beta   & 3.68  & .036 & * \\
        EV\_007 & \gls{FAI}       & 4.73  & .046 & ** \\
        EV\_008 & Frontal beta    & 11.62 & .110 & *** \\
        \hline
    \end{tabular}
    \par\vspace{4pt}
    \noindent\parbox{\linewidth}{\scriptsize Only effects with $p < .05$ are listed. All $df_1 = 2$; $df_2$ varies by modality ($\approx$160--200 for EEG). Non-significant participants and bands omitted for brevity; see Table~\ref{tab:pilot_anova} for the complete breakdown.}
\end{table}

\subsubsection{Prefrontal Hemodynamic Responses: fNIRS}

\gls{fNIRS} measured cerebral hemodynamic changes (\gls{HbO2} and \gls{HbR} concentration changes) in prefrontal cortex during emotional processing. Technical issues (excessive motion artifacts, poor optode coupling) compromised data quality for three of five participants — a known challenge in multimodal setups \parencite{scholkmannReviewContinuousWave2014}. Interpretable hemodynamic data were obtained only for EV\_002 and EV\_003.

For these two participants, condition-dependent prefrontal activation patterns were observed, and \gls{fNIRS}-derived frontal asymmetry indices complemented the \gls{EEG}-based \gls{FAI}, providing convergent hemodynamic evidence for lateralized prefrontal engagement. All \gls{fNIRS} \gls{ANOVA} effects were non-significant across both participants (all $p > .06$; see Table~\ref{tab:pilot_anova}). The key data characteristics are summarised in Table~\ref{tab:ev1_fnirs}; hemodynamic traces are available in \texttt{EV\_results.html}.

\begin{table}[htbp]
    \centering
    \small
    \caption{EV1 fNIRS Data Summary (valid participants only)}
    \label{tab:ev1_fnirs}
    \begin{tabular}{lp{4cm}p{5.5cm}}
        \hline
        \textbf{Participant} & \textbf{Data quality} & \textbf{Observed pattern} \\
        \hline
        EV\_002 & Valid (2/5 participants) & Condition-dependent \gls{HbO2} modulation in prefrontal \glspl{ROI}; \gls{FAI} trend ($p = .064$, Left PFC) \\
        EV\_003 & Valid & Prefrontal activation present; all \gls{ANOVA} effects ns \\
        EV\_004 & Excluded — motion artifacts & --- \\
        EV\_007 & Excluded — poor optode coupling & --- \\
        EV\_008 & Excluded — poor optode coupling & --- \\
        \hline
    \end{tabular}
    \par\vspace{4pt}
    \noindent\parbox{\linewidth}{\scriptsize ns = not significant ($p > .05$). Hemodynamic time courses and ROI bar charts for EV\_002 and EV\_003 available in \texttt{EV\_results.html}.}
\end{table}

\subsubsection{Within-Participant Statistical Summary}
\label{sec:pilot_anova_summary}

Table~\ref{tab:pilot_anova} summarizes the within-participant one-way \gls{ANOVA} results across all five participants for each modality stream. The factor was valence condition (NEG, NEU, POS; df1\,=\,2). Degrees of freedom (df2) varied by modality due to differences in the number of valid sliding-window epochs per participant. All results are uncorrected for multiple comparisons and should be treated as exploratory.

\begin{table}[htbp]
    \centering
    \small
    \setlength{\tabcolsep}{5pt}
    \caption{Within-Participant One-Way \gls{ANOVA} Results: Study 1 (EmotiView Pilot, $N = 5$)}
    \label{tab:pilot_anova}
    \begin{tabular}{llrrrrrrrr}
        \hline
        & & \multicolumn{2}{c}{\textbf{EEG Frontal}} & \multicolumn{2}{c}{\textbf{EEG Parietal}} & & & & \\
        \cline{3-4}\cline{5-6}
        \textbf{Participant} & & \textbf{theta} & \textbf{beta} & \textbf{alpha} & \textbf{beta} & \textbf{FAI} & \textbf{HRV} & \textbf{EDA} & \textbf{fNIRS} \\
        \hline
        EV\_002 & $F$ & 4.00 & 1.21 & 0.08 & 0.28 & 0.39 & 27.80 & 0.33 & 5.03 / 2.07 \\
               & $p$ & .020 & .302 & .923 & .754 & .679 & $<$.001 & .730 & .064 / .221 \\
               & $\eta^2$ & .048 & .015 & .001 & .003 & .004 & .227 & .099 & .668 / .453 \\
               & sig & * & ns & ns & ns & ns & *** & ns & ns / ns \\
        \hline
        EV\_003 & $F$ & 34.92 & 2.03 & 0.46 & 0.94 & 4.50 & 11.02 & 1.35 & 0.99 / 0.11 \\
               & $p$ & $<$.001 & .134 & .633 & .392 & .012 & $<$.001 & .329 & .434 / .897 \\
               & $\eta^2$ & .297 & .024 & .005 & .010 & .047 & .104 & .310 & .284 / .043 \\
               & sig & *** & ns & ns & ns & * & *** & ns & ns / ns \\
        \hline
        EV\_004 & $F$ & 0.75 & 1.02 & 1.44 & 0.37 & 1.75 & 2.49 & 0.39 & 1.97 / 0.44 \\
               & $p$ & .472 & .362 & .241 & .691 & .177 & .086 & .695 & .220 / .665 \\
               & $\eta^2$ & .008 & .010 & .015 & .004 & .018 & .028 & .114 & .396 / .127 \\
               & sig & ns & ns & ns & ns & ns & ns & ns & ns / ns \\
        \hline
        EV\_007 & $F$ & 0.71 & 0.81 & 3.13 & 3.68 & 4.73 & 3.87 & 0.67 & 1.81 / 2.44 \\
               & $p$ & .493 & .446 & .046 & .027 & .010 & .023 & .548 & .243 / .168 \\
               & $\eta^2$ & .008 & .009 & .031 & .036 & .046 & .039 & .182 & .376 / .448 \\
               & sig & ns & ns & * & * & ** & * & ns & ns / ns \\
        \hline
        EV\_008 & $F$ & 1.43 & 11.62 & 0.71 & 0.75 & 0.52 & 18.71 & 0.69 & 0.09 / 1.00 \\
               & $p$ & .242 & $<$.001 & .491 & .475 & .593 & $<$.001 & .538 & .913 / .423 \\
               & $\eta^2$ & .015 & .110 & .007 & .008 & .005 & .164 & .187 & .030 / .249 \\
               & sig & ns & *** & ns & ns & ns & *** & ns & ns / ns \\
        \hline
        \multicolumn{10}{l}{\rule{0pt}{8pt}} \\[-6pt]
        \hline
    \end{tabular}
    \par\vspace{4pt}
    \noindent\parbox{\linewidth}{\scriptsize * $p < 0.05$; ** $p < 0.01$; *** $p < 0.001$. All df1\,=\,2. fNIRS columns: Left\,PFC / Right\,PFC. df2 varies by modality (EEG: $\approx$160--200; HRV/EDA: $\approx$6--9; fNIRS: 5--6). \textbf{Note:} fNIRS $\eta^2$ values are artefactually inflated due to df2\,$\leq$\,6 ($\approx$2--3 observations per condition cell) and should not be interpreted as meaningful effect-magnitude estimates.}
\end{table}

\noindent Notable patterns across participants: \gls{HRV} showed the most consistent condition effects, reaching significance in 4 of 5 participants (EV\_002, EV\_003, EV\_007, EV\_008), with large effect sizes in two (EV\_002: $\eta^2 = 0.23$; EV\_008: $\eta^2 = 0.16$). Frontal theta power was significant in 2 of 5 participants (EV\_002: $\eta^2 = 0.05$; EV\_003: $\eta^2 = 0.30$), with the latter representing an unusually large effect. \gls{FAI} reached significance in EV\_003 and EV\_007. \gls{EDA} and \gls{fNIRS} were non-significant across all participants. EV\_004 showed no significant effects in any modality.

\vspace{1em}

% Summary of Pilot Findings moved to Discussion section
% (interpretation and evaluation of pilot results belongs there)
% ===================================================================
\subsection{Study 2: Secondary Validation Using the DEAP Dataset}
\label{sec:results_ev2}
% ===================================================================

\subsubsection{Pipeline Completion and Data Coverage}

The EV2 analysis pipeline executed successfully for all 32 DEAP participants (DEAP\_01 through DEAP\_32), each completing 15 terminal output modules within a single Nextflow session. Processing durations ranged from approximately 1 s to 452 s per participant due to Nextflow's parallel scheduling, with the full cohort finalizing by 12:17 (UTC+2) on 2026-06-08. Group-level (L2) aggregation produced an \texttt{l2\_anova\_summary.parquet} containing pooled one-way \gls{ANOVA} results for five modality streams (EDA, \gls{FAI}, EEG Frontal band power, EEG Parietal band power, \gls{PRV}). All 32 participant result trees were committed to version control, ensuring full reproducibility of downstream analyses and visualizations.

\subsubsection{Affective Rating Distribution}

The DEAP rating distributions reflect the broad stimulus design of the dataset, spanning the full 1--9 continuum in both valence and arousal. Across all 32 participants and 40 trials ($N = 1{,}280$ trial-level observations), valence ratings covered the full scale (population-level range approximately 0--9), while arousal ratings similarly spanned low-to-high activation states. As a representative example, participant DEAP\_01 provided valence ratings with $M = 3.12$ ($SD = 2.16$; range $[0.00, 8.27]$) and arousal ratings with $M = 3.76$ ($SD = 2.33$; range $[1.00, 8.15]$), indicating a broad spread across the affective plane. This wide distribution is characteristic of DEAP's 40-clip design, which intentionally samples diverse valence--arousal combinations \parencite{koelstraDeapDatabaseEmotion2012}, and is necessary to support the trial-binning procedure (threshold = 5.0, neutral band $\pm 1.0$; see Section~\ref{sec:ev2_binning}).

The trial-binning procedure successfully generated up to six condition bins (arousal\_high, arousal\_low, valence\_high, valence\_low, valence\_neutral, and, where supported by the rating distribution, valence\_high\_arousal\_high) for each participant. The number of trials per bin ranged from approximately 6 to 18, reflecting individual differences in how each participant's ratings mapped onto the threshold boundaries. Participants whose ratings clustered near the neutral midpoint (4--6) contributed fewer trials to high/low extreme bins, which directly affected the statistical precision of within-participant bootstrap estimates and \gls{ANOVA} tests.

\subsubsection{Frontal Alpha Asymmetry}

\paragraph{Within-participant \gls{FAI} profiles.}
Per-trial \gls{FAI} ($\ln P_{\text{F4}} - \ln P_{\text{F3}}$) was computed for all 40 trials per participant. Participant DEAP\_01 exhibited a consistent right-hemisphere alpha power dominance, yielding a uniformly negative \gls{FAI} across all 40 trials ($M = -0.323$, $SD = 0.161$; range $[-0.897, -0.063]$; 100\% of trials with $\text{FAI} < 0$). This pattern, indicative of sustained relative right-frontal activation irrespective of valence condition, was present in a subset of participants and likely reflects stable inter-individual differences in hemispheric asymmetry \parencite{allenIssuesAssumptionsRoad2004}. Other participants showed more mixed trial-level distributions, with the direction of \gls{FAI} varying across the 40 clips. Across the sample, \gls{FAI} values covered the range from strongly left-dominant (positive values, approach motivation) to strongly right-dominant (negative values, withdrawal), consistent with the expected individual variability reported in the \gls{FAI} literature \parencite{smithPrefrontalAsymmetryMeta2017}.

\paragraph{Bootstrap confidence intervals across condition bins.}
Bootstrap 95\% percentile confidence intervals (CIs; 10{,}000 iterations) for \gls{FAI} by condition bin are reported for the representative participant DEAP\_01 in Table~\ref{tab:deap01_fai_bs}. Interactive bar-chart visualizations for all 32 participants are available in the pipeline HTML archive (\texttt{EV2\_results.html}). The bootstrap analysis yielded the following condition-level estimates:

\begin{table}[htbp]
    \centering
    \caption{Bootstrap \gls{FAI} Estimates by Condition Bin (DEAP\_01, representative participant)}
    \label{tab:deap01_fai_bs}
    \begin{tabular}{lrrr}
        \hline
        \textbf{Condition} & \textbf{Mean FAI} & \textbf{95\% CI lower} & \textbf{95\% CI upper} \\
        \hline
        Arousal High  & $-0.232$ & $-0.297$ & $-0.168$ \\
        Arousal Low   & $-0.360$ & $-0.426$ & $-0.303$ \\
        Valence High  & $-0.304$ & $-0.358$ & $-0.248$ \\
        Valence Low   & $-0.329$ & $-0.394$ & $-0.271$ \\
        Valence Neutral & $-0.304$ & $-0.440$ & $-0.150$ \\
        \hline
    \end{tabular}
\end{table}

\noindent A notable pattern is observable in the arousal dimension: the arousal-high bin yielded a higher (less negative) \gls{FAI} than the arousal-low bin (difference $\approx 0.13$ log-units), with non-overlapping bootstrap CIs. The valence-based bins showed largely overlapping CIs. The wider CI for valence\_neutral reflects the smaller number of trials allocated to this bin.

\paragraph{Within-participant one-way \gls{ANOVA}.}
The within-participant \gls{ANOVA} for \gls{FAI} across condition bins was non-significant for DEAP\_01 ($F(5, 74) = 1.03$, $p = 0.407$, $\eta^2 = 0.065$) and similarly non-significant for the majority of participants. This is consistent with the bootstrap pattern: the arousal modulation observed in the bootstrap plot does not reach conventional significance levels at $n \approx 10$--18 trials per bin, and effect sizes remained in the small-to-moderate range.

% Figure: bootstrap FAI by condition (rendered in HTML archive EV2_results.html, not available as standalone PDF)
% \label{fig:deap01_fai_bs}

\subsubsection{EEG Region-of-Interest Band Power}

\paragraph{Within-participant \gls{ANOVA}.}
The one-way \gls{ANOVA} applied to the sliding-window EEG \gls{ROI} band power (30 s windows, 10 s step, $n \approx 4$ windows/trial, $\approx 160$ windows/participant) revealed condition-dependent band power modulation in a substantial proportion of participants. Within the four representative participants whose data were inspected (DEAP\_01, DEAP\_05, DEAP\_10, DEAP\_20), significant condition effects were observed consistently in the theta band:

\begin{itemize}
    \item DEAP\_01: Frontal theta $F(5, 234) = 4.75$, $p < 0.001$, $\eta^2 = 0.092$; Parietal theta $F(5, 234) = 7.04$, $p < 0.001$, $\eta^2 = 0.131$.
    \item DEAP\_05: Frontal theta $F(5, 234) = 4.50$, $p = 0.001$, $\eta^2 = 0.088$; Parietal theta $F(5, 234) = 3.89$, $p = 0.002$, $\eta^2 = 0.077$.
    \item DEAP\_10: Frontal theta $F(5, 234) = 4.70$, $p < 0.001$, $\eta^2 = 0.091$; Frontal beta $F(5, 234) = 4.00$, $p = 0.002$, $\eta^2 = 0.079$.
    \item DEAP\_20: Parietal beta $F(5, 234) = 9.48$, $p < 0.001$, $\eta^2 = 0.168$ (the largest individual effect observed).
\end{itemize}

The alpha band showed significant effects in several participants (e.g., DEAP\_05 frontal alpha $F(5, 234) = 3.65$, $p = 0.003$; DEAP\_10 parietal alpha $F(5, 234) = 4.14$, $p = 0.001$), while beta power was significant in participants with higher trial-count bins. The pattern suggests that individual participants show genuine condition-level modulation of spectral power, with the theta band showing the most consistent effects across participants.

\paragraph{OLS valence contrasts.}
For the valence contrast (vH $-$ vL), participant DEAP\_01 showed small and non-significant differences across all frontal bands: alpha ($\Delta = -0.005$, $t = -0.042$), beta ($\Delta = -0.025$, $t = -0.528$), theta ($\Delta = -0.142$, $t = -0.993$). In the parietal \gls{ROI}, parietal alpha showed a small positive vH $-$ vL contrast ($\Delta = +0.108$, $t = +0.746$), while beta and theta were near zero. These effect-level contrasts indicate that the valence-based partitioning did not produce strong power differences for this participant, consistent with the \gls{FAI} bootstrap pattern described above.

\subsubsection{Electrodermal Activity}

Per-trial \gls{EDA} amplitude (mean conductance level, µS) exhibited wide within-participant variability consistent with individual differences in electrodermal responsivity. Bootstrap CIs for DEAP\_01 revealed that the arousal-high bin yielded a numerically higher EDA amplitude ($M = 8{,}822\,\mu\text{S}$, 95\% CI $[2{,}865,\, 14{,}696]$) compared to the arousal-low bin ($M = 5{,}268\,\mu\text{S}$, 95\% CI $[-4{,}817,\, 13{,}675]$), although the CIs overlap substantially. The valence-neutral bin showed the highest mean EDA amplitude ($M = 16{,}106\,\mu\text{S}$, 95\% CI $[9{,}947,\, 19{,}232]$), which likely reflects baseline drift rather than a valence-specific response, since no tonic/phasic decomposition was applied to isolate stimulus-evoked \glspl{SCR} from slow skin conductance level trends. The within-participant \gls{ANOVA} for EDA amplitude across bins was non-significant for DEAP\_01 ($F(5, 74) = 0.20$, $p = 0.961$, $\eta^2 = 0.013$).



\subsubsection{Pulse Rate Variability}

The per-trial \gls{PRV}--\gls{RMSSD} (derived from BVP peak intervals at 128 Hz, minimum inter-peak distance 0.4 s) showed the expected inverse relationship with arousal for DEAP\_01: the arousal-high bin yielded $M = 159.3\,\text{ms}$ (95\% CI $[149.5,\, 169.6]$) compared to $M = 163.3\,\text{ms}$ (95\% CI $[158.3,\, 168.4]$) for arousal-low, consistent with the theoretical prediction that higher sympathetic activation reduces vagal-mediated heart rate variability \parencite{malikHeartRateVariability1996, shafferOverviewHeartRate2017}. The valence-neutral bin showed the lowest \gls{RMSSD} ($M = 146.8\,\text{ms}$, 95\% CI $[144.7,\, 149.6]$).

\noindent\textbf{Data quality note.} The \gls{RMSSD} values observed here ($\approx 146$--$166$\,ms) substantially exceed typical \gls{HRV} values for resting adults regardless of recording modality ($\approx 30$--$50$\,ms for ECG; pulse-to-pulse intervals from photoplethysmography closely approximate ECG-based R-R interval variability when measured under low-noise conditions \parencite{allenPhotoplethysmographyItsApplication2007}). The elevated values are most likely attributable to false-positive peak detections at secondary peaks in the dicrotic notch of the unprocessed BVP waveform at 128\,Hz, yielding artificially short inter-peak intervals that inflate \gls{RMSSD} \parencite{pengrosaPioneerPhysiologicalOptics2024}. The minimum inter-peak distance parameter (0.4\,s) mitigates but does not eliminate this issue.

The within-participant \gls{ANOVA} for \gls{PRV}--\gls{RMSSD} was non-significant for DEAP\_01 ($F(5, 74) = 1.03$, $p = 0.407$).

\subsubsection{Group-Level Linear Mixed Model}

The L2 group-level analysis used a Linear Mixed Model (LMM) with condition bin as the fixed effect and participant random intercept, as described in the Methods (Section~\ref{sec:ev2_methods}). This approach properly partitions the variance into within-condition fixed effects and between-participant random variation, avoiding the pseudoreplication inherent in pooling all trials across participants into a single ANOVA. The full pipeline was re-run after the model change; results are summarised in Table~\ref{tab:l2_lmm}.

\begin{table}[htbp]
    \centering
    \caption{Group-Level (L2) LMM Results Across Modality Streams ($N = 32$ participants; random intercept per participant)}
    \label{tab:l2_lmm}
    \small
    \begin{tabular}{llrrrrrrl}
        \hline
        \textbf{Modality} & \textbf{DV} & $\boldsymbol{\chi^2}$ & \textbf{df} & $\boldsymbol{p}$ & $\boldsymbol{R^2_m}$ & $\boldsymbol{R^2_c}$ & \textbf{ICC} & \textbf{sig} \\
        \hline
        EDA          & amplitude & $\infty$\textsuperscript{\dag} & 5 & $<.001$ & --- & --- & --- & *** \\
        \gls{FAI}    & fai       & 12.40 & 5 & .030    & .500 & .979 & .958 & *   \\
        \gls{PRV}    & RMSSD     & $\infty$\textsuperscript{\dag} & 5 & $<.001$ & --- & --- & --- & *** \\
        EEG Frontal  & alpha     & 16.78 & 5 & .005    & .499 & .888 & .778 & **  \\
        EEG Frontal  & beta      & 31.67 & 5 & $<.001$ & .499 & .916 & .832 & *** \\
        EEG Frontal  & theta     & 34.80 & 5 & $<.001$ & .499 & .857 & .715 & *** \\
        EEG Parietal & alpha     &  5.18 & 5 & .394    & .498 & .795 & .593 & ns  \\
        EEG Parietal & beta      & 15.79 & 5 & .008    & .499 & .890 & .781 & **  \\
        EEG Parietal & theta     & 13.58 & 5 & .019    & .498 & .792 & .586 & *   \\
        \hline
    \end{tabular}
    \par\vspace{4pt}
    \noindent\parbox{\linewidth}{\scriptsize $\chi^2$ = LRT chi-squared statistic (full vs.\ null model, ML estimation); $R^2_m$ = marginal $R^2$ (fixed effects only); $R^2_c$ = conditional $R^2$ (fixed + random); ICC = intraclass correlation (between-participant variance share). $n_\text{obs}$ = 2{,}560 for EDA/\gls{FAI}/\gls{PRV}; 7{,}680 for EEG streams ($N = 32$ participants). * $p < .05$; ** $p < .01$; *** $p < .001$. \textsuperscript{\dag}~Degenerate LRT arising from a singular random-effects solution (random intercept variance $\rightarrow 0$); $R^2$ and ICC are not computable; omnibus significance should be interpreted with caution (see Discussion).}
\end{table}

The ICC values will quantify the proportion of total variance attributable to stable between-participant differences; a high ICC indicates that participant identity is a strong predictor of baseline physiology, reinforcing the necessity of the random-effects structure. The marginal $R^2_m$ will reflect the condition fixed effect after controlling for this participant-level variance, providing an unconfounded estimate of the effect of valence and arousal binning on each modality stream.

\subsubsection{Cross-modal Consistency}

The multimodal cross-correlation analysis (computed per-participant across the trial-level feature vectors) revealed strong within-modality EEG band correlations for DEAP\_01: alpha--beta Spearman $r = 0.588$ ($p < 0.001$), alpha--theta $r = 0.717$ ($p < 0.001$), consistent with the joint modulation of multiple EEG frequency bands by sustained spectral states during long stimulus epochs. The \gls{FAI} measure (reflecting anterior asymmetry) showed weaker correlation with the global band power metrics, as expected since it captures a lateralization ratio rather than absolute power. The BVP-derived interval feature showed moderate correlation with EEG alpha ($r = 0.449$, $p < 0.01$).

% Summary of Study 2 Findings moved to Discussion section
